\section{The truncated approximation of a filtered term}
\label{sec:hastings_truncated_dynamics}

% Hastings, arXiv:cond-mat/0508554v3, proj3new.tex, eqs. mid, odef, hpdef, apb
% (lines 285--320), proved at lines 1458--1533.

For a term $H_i$ of a local Hamiltonian $H=\sum_jH_j$, \cite{Hastings2006Solving}
approximates the filtered term $\tilde H_i^0$ by the filter of $H_i$ under the
truncated Hamiltonian $H_{\rm loc}=\sum_{j:\,d(i,j)\le l_{proj}-R}H_j$,
\begin{align*}
    M_i=\frac{\Delta E}{\sqrt{2\pi q}}\int dt\,e^{-(t\Delta E)^2/(2q)}\,
    e^{iH_{\rm loc}t}H_ie^{-iH_{\rm loc}t},
\end{align*}
\cite[eqs.~(mid), (odef), (hpdef), lines 285--303]{Hastings2006Solving}. Here
the truncation keeps any family $s$ of terms acting inside a set $K$ of sites
containing the support $S$ of the filtered operator $W$, and the shell consists
of the omitted terms whose supports meet $K$. The time integral of the
Lieb--Robinson bound is bounded by $T$ times its largest value, the constants
are explicit, and the Gaussian weight of eq.~(mid) is read as that of
eq.~(thd) \cite{gap:hastings05_truncated_dynamics}.

\begin{theorem}[Support of the truncated filter]
    \label{thm:hastings_truncated_support}
    \lean{TNLean.Lattice.hermitianUnitaryPath_conj_mem_supportedOperators,
        TNLean.Lattice.gaussianFilteredTerm_mem_supportedOperators}
    \leanok
    \uses{def:hastings_gaussian_filtered_term}
    If $H$ and $W$ act on a set $K$ of sites, then so do $e^{iHt}We^{-iHt}$ for
    every $t$ and the filtered operator $\widetilde W$. In particular $M_i$ acts on
    the sites within distance $l_{proj}$ of $i$
    \cite[lines 302--304]{Hastings2006Solving}.
\end{theorem}

\begin{proof}
    \leanok
    The operators acting on $K$ form a closed subalgebra, so they contain the
    exponential $e^{iHt}$ and the products $e^{iHt}We^{-iHt}$. The filter is
    the average of these products against a probability measure, and a
    closed subspace contains such averages.
\end{proof}

\begin{theorem}[Comparison of two evolutions]
    \label{thm:hastings_truncated_evolution_comparison}
    \lean{TNLean.Lattice.norm_conj_sub_conj_le}
    \leanok
    Let $H$ be Hermitian and $H'$ any matrix, and write
    $\tau'_u(W)=e^{iH'u}We^{-iH'u}$. If
    $\|[H-H',\tau'_u(W)]\|\le\varepsilon$ for all $u$ between $0$ and $t$, then
    \begin{align*}
        \|e^{iHt}We^{-iHt}-e^{iH't}We^{-iH't}\|\le\varepsilon\,|t|.
    \end{align*}
    This is the comparison of $H_i(t)$ and $H_i^{\rm trunc}(t)$ in
    \cite[lines 1469--1482]{Hastings2006Solving}.
\end{theorem}

\begin{proof}
    \leanok
    Put $f(s)=e^{iHs}\tau'_{t-s}(W)e^{-iHs}$, so that $f(0)=\tau'_t(W)$ and
    $f(t)=e^{iHt}We^{-iHt}$. Differentiating,
    $f'(s)=e^{iHs}\,i[H-H',\tau'_{t-s}(W)]\,e^{-iHs}$, whose norm is at most
    $\varepsilon$ because $e^{iHs}$ is unitary. The mean value inequality on
    the interval between $0$ and $t$ gives the bound.
\end{proof}

\begin{theorem}[Filtering a comparison of evolutions]
    \label{thm:hastings_truncated_filter_comparison}
    \lean{TNLean.Lattice.norm_gaussianFilteredTerm_sub_le}
    \leanok
    \uses{def:hastings_gaussian_filtered_term}
    Let $H$ and $H'$ be Hermitian, $q\ge0$, $T\ge0$ and $\varepsilon\ge0$. If
    $\|e^{iHt}We^{-iHt}-e^{iH't}We^{-iH't}\|\le\varepsilon$ for all $|t|\le T$,
    then the filters of $W$ under $H$ and under $H'$ differ in norm by at most
    $\varepsilon+2\|W\|e^{-(T\Delta E)^2/(2q)}$
    \cite[eq.~(dmi) and lines 1524--1533]{Hastings2006Solving}.
\end{theorem}

\begin{proof}
    \leanok
    \uses{thm:hastings_filtered_locality}
    The difference of the filters is the Gaussian average of the difference of
    the evolutions. Split the average at $|t|=T$: inside, the integrand is at
    most $\varepsilon$; outside, it is at most $2\|W\|$, and the Gaussian mass
    of the complement of $[-T,T]$ is at most $e^{-(T\Delta E)^2/(2q)}$.
\end{proof}

\begin{theorem}[Lieb--Robinson bound against an operator on several sites]
    \label{thm:hastings_truncated_multisite_lieb_robinson}
    \lean{TNLean.Lattice.norm_commutator_heisenbergEvolution_le_graphDistance_finset}
    \leanok
    \uses{thm:qc_graph_lieb_robinson_finset}
    Under the finite-range interaction assumptions of
    Theorem~\ref{thm:qc_graph_lieb_robinson_finset}, let $A$ act on a finite set
    $S$ and $B$ on a finite set $Y$, with $m\le d_G(x,y)$ for all $x\in S$ and
    $y\in Y$. For every $\mu\ge0$,
    \begin{align*}
        \|[B,\tau_t(A)]\|\le2|Y|\,C\,|S|\,\|A\|\,\|B\|\,e^{v|t|-\mu m},
    \end{align*}
    with $C=2+4|b_0|e^{2\mu R}$ and $v=2v_R|b_0|e^{2\mu R}+1$.
\end{theorem}

\begin{proof}
    \leanok
    \uses{thm:qc_graph_lieb_robinson_finset, lem:qc_site_expectation_commutator}
    Let $E$ be the normalized partial trace over $Y$. Then $E(\tau_t(A))$ acts
    outside $Y$ and commutes with $B$, so
    $\|[B,\tau_t(A)]\|\le2\|B\|\,\|\tau_t(A)-E(\tau_t(A))\|$. By
    Lemma~\ref{lem:qc_site_expectation_commutator}, the last norm is at most the
    sum over $y\in Y$ of the single-site commutator bounds of
    Theorem~\ref{thm:qc_graph_lieb_robinson_finset}, each at most
    $C|S|\,\|A\|e^{v|t|-\mu m}$.
\end{proof}

\begin{theorem}[The omitted terms against the truncated evolution]
    \label{thm:hastings_truncated_omitted_commutator}
    \lean{TNLean.Lattice.norm_commutator_sub_truncated_le}
    \leanok
    \uses{thm:hastings_truncated_multisite_lieb_robinson}
    Let $W$ act on $S\subseteq K$, let $s$ be a family of terms acting inside
    $K$, put $H_s=\sum_{k\in s}h_k$, and let $m$ bound from below the graph
    distance from $S$ to the support of every omitted term meeting $K$. Then
    \begin{align*}
        \|[H-H_s,e^{iH_su}We^{-iH_su}]\|
        \le2v_RC\,|S|\,\|W\|\Bigl(\sum_k\|h_k\|\Bigr)e^{v|u|-\mu m},
    \end{align*}
    the sum running over the omitted terms meeting $K$
    \cite[lines 1469--1505]{Hastings2006Solving}.
\end{theorem}

\begin{proof}
    \leanok
    \uses{thm:hastings_truncated_support,
        thm:hastings_truncated_multisite_lieb_robinson}
    The difference $H-H_s$ is the sum of the omitted terms. The evolution under
    $H_s$ acts on $K$, so an omitted term whose support is disjoint from $K$
    commutes with it. For an omitted term meeting $K$, apply
    Theorem~\ref{thm:hastings_truncated_multisite_lieb_robinson} to the family
    in which the terms outside $s$ are replaced by $0$, which generates the
    evolution under $H_s$ and satisfies the same interaction bounds, and use
    $|X_k|\le v_R$.
\end{proof}

\begin{theorem}[Hastings' truncated approximation]
    \label{thm:hastings_truncated_approximation}
    \lean{TNLean.Lattice.norm_gaussianFilteredTerm_sub_truncated_le_sum,
        TNLean.Lattice.norm_gaussianFilteredTerm_sub_truncated_le,
        TNLean.Lattice.truncated_approx_gaussianFilteredTerm}
    \leanok
    \uses{def:hastings_gaussian_filtered_term}
    In the setting of Theorem~\ref{thm:hastings_truncated_omitted_commutator},
    let every term have norm at most $J$ and let $N$ be the number of omitted
    terms meeting $K$. Let $M$ be the filter of $W$ under $H_s$. Then $M$ acts
    on $K$, $\|M\|\le\|W\|$, $M$ is Hermitian when $W$ is, and for every
    $q\ge0$, $T\ge0$ and $\mu\ge0$,
    \begin{align*}
        \|\widetilde W-M\|\le2v_RC\,|S|\,\|W\|\,J\,N\,T\,e^{vT-\mu m}
        +2\|W\|e^{-(T\Delta E)^2/(2q)}.
    \end{align*}
    With $W=H_i$, $K$ the sites within distance $l_{proj}$ of $i$ and $s$ the
    terms with $d(i,j)\le l_{proj}-R$, the operator $M$ is $M_i$, and this is
    the analogue of \cite[eq.~(apb), lines 300--320]{Hastings2006Solving}
    with explicit constants and $q$, $T$ free. The comparison of the shell
    count with $N(l_{proj})$ for this choice is recorded in
    \cite{gap:hastings05_truncated_dynamics}.
\end{theorem}

\begin{proof}
    \leanok
    \uses{thm:hastings_truncated_support,
        thm:hastings_truncated_evolution_comparison,
        thm:hastings_truncated_filter_comparison,
        thm:hastings_truncated_omitted_commutator,
        thm:hastings_gaussian_filtered_spectral}
    For $|t|\le T$ and $u$ between $0$ and $t$,
    Theorem~\ref{thm:hastings_truncated_omitted_commutator} bounds the
    commutator by $\varepsilon=2v_RC|S|\,\|W\|JN\,e^{vT-\mu m}$, so
    Theorem~\ref{thm:hastings_truncated_evolution_comparison} bounds the
    difference of the evolutions by $\varepsilon T$, and
    Theorem~\ref{thm:hastings_truncated_filter_comparison} gives the estimate.
    The support is Theorem~\ref{thm:hastings_truncated_support}, the norm bound
    is that of Theorem~\ref{thm:hastings_gaussian_filtered_spectral}, and the
    adjoint of the filter of $W$ is the filter of $W^\dagger$.
\end{proof}

\begin{theorem}[The truncated approximation at the time window $T=q/\Delta E$]
    \label{thm:hastings_truncated_approximation_time}
    \lean{TNLean.Lattice.exp_neg_sq_div_eq_of_time_eq,
        TNLean.Lattice.norm_gaussianFilteredTerm_sub_truncated_le_of_time}
    \leanok
    \uses{thm:hastings_truncated_approximation}
    In the setting of Theorem~\ref{thm:hastings_truncated_approximation}, let
    $q\ge0$ and $\Delta E>0$. At $T=q/\Delta E$ the Gaussian tail
    $e^{-(T\Delta E)^2/(2q)}$ equals $e^{-q/2}$, so
    \begin{align*}
        \|\widetilde W-M\|\le2v_RC\,|S|\,\|W\|\,J\,N\,\frac{q}{\Delta E}\,
        e^{vq/\Delta E-\mu m}+2\|W\|e^{-q/2}.
    \end{align*}
    For $q=c_1l_{proj}\Delta E$ this is the time $T=c_1l_{proj}$ of
    \cite[lines 337--349]{Hastings2006Solving}.
\end{theorem}

\begin{proof}
    \leanok
    \uses{thm:hastings_truncated_approximation}
    Put $T=q/\Delta E$ in Theorem~\ref{thm:hastings_truncated_approximation};
    then $(T\Delta E)^2/(2q)=q/2$ for $q>0$, and both sides of the tail
    identity are $1$ for $q=0$.
\end{proof}

\begin{theorem}[Ground vectors under the truncated approximation]
    \label{thm:hastings_truncated_facld}
    \lean{TNLean.Lattice.norm_sub_starProjection_truncated_apply_le,
        TNLean.Lattice.norm_truncated_apply_sub_inner_smul_le}
    \leanok
    \uses{thm:hastings_truncated_approximation_time, thm:hastings05_facld}
    In the setting of Theorem~\ref{thm:hastings_truncated_approximation}, let
    $q\ge0$ and $\Delta E>0$, and let every eigenvalue of $H=\sum_kh_k$ equal
    $E_0$ or be at least $E_0+\Delta E$. Write
    \begin{align*}
        \delta=3\|W\|e^{-q/2}+2v_RC\,|S|\,\|W\|\,J\,N\,\frac{q}{\Delta E}\,
        e^{vq/\Delta E-\mu m}.
    \end{align*}
    For every ground vector $\psi$, the component of $M\psi$ orthogonal to the
    ground space has norm at most $\delta\|\psi\|$. If $\Omega$ is a unique
    ground state with gap $\Delta E$, then
    $\|M\Omega-\langle\Omega,M\Omega\rangle\Omega\|\le\delta$. With $W=H_i$
    and the choices of Theorem~\ref{thm:hastings_truncated_approximation},
    these are eq.~(facld) of \cite[lines 322--348]{Hastings2006Solving} and
    the bound for one ground state at lines 375--387, with explicit
    constants and $q$ free.
\end{theorem}

\begin{proof}
    \leanok
    \uses{thm:hastings_truncated_approximation_time, thm:hastings05_facld}
    Theorem~\ref{thm:hastings05_facld} bounds both quantities by
    $\|W\|e^{-q/2}+\|\widetilde W-M\|$ (times $\|\psi\|$ for the first), and
    Theorem~\ref{thm:hastings_truncated_approximation_time} bounds
    $\|\widetilde W-M\|$.
\end{proof}

\begin{theorem}[The sum of the truncated approximations]
    \label{thm:hastings_truncated_sum}
    \lean{TNLean.Lattice.norm_sum_truncated_sub_le,
        TNLean.Lattice.norm_sum_truncated_sub_le_card}
    \leanok
    \uses{thm:hastings_truncated_approximation_time,
        thm:hastings05_sum_approximation}
    In the setting of Theorem~\ref{thm:hastings_truncated_approximation}, let
    $q\ge0$, $\Delta E>0$ and $\mu\ge0$, and let $H=\sum_ih_i$ be Hermitian. For each
    term $h_i$, acting on $X_i$, let $s_i$ be a set of terms acting inside a
    set of sites $K_i\supseteq X_i$, such that the omitted terms meeting $K_i$
    are at distance at least $m$ from $X_i$. Let $N_i$ be their number, and let
    $M_i$ be the filter of $h_i$ under $\sum_{k\in s_i}h_k$. Then
    \begin{align*}
        \Bigl\|\sum_iM_i-H\Bigr\|\le\sum_i\Bigl(2v_RC\,|X_i|\,\|h_i\|\,J\,N_i\,
        \frac{q}{\Delta E}\,e^{vq/\Delta E-\mu m}+2\|h_i\|e^{-q/2}\Bigr).
    \end{align*}
    If $N_i\le N_0$ for every $i$, the right-hand side is at most the number
    of terms times
    $2v_R^2C\,J^2N_0\,(q/\Delta E)\,e^{vq/\Delta E-\mu m}+2Je^{-q/2}$. With
    one term per site and the choices of
    Theorem~\ref{thm:hastings_truncated_approximation}, this is the bound on
    the difference between $H$ and $M=\sum_iM_i$ by the number of sites times
    $\|\widetilde H_i^0-M_i\|$ in \cite[lines 343--374]{Hastings2006Solving},
    with explicit constants and $q$ free.
\end{theorem}

\begin{proof}
    \leanok
    \uses{thm:hastings_truncated_approximation_time,
        thm:hastings05_sum_approximation}
    Theorem~\ref{thm:hastings05_sum_approximation} bounds
    $\|\sum_iM_i-H\|$ by $\sum_i\|\widetilde h_i-M_i\|$, and
    Theorem~\ref{thm:hastings_truncated_approximation_time} with $W=h_i$ and
    $S=X_i$ bounds each summand. For the uniform form use $|X_i|\le v_R$,
    $\|h_i\|\le J$ and $N_i\le N_0$.
\end{proof}

\begin{remark}
    Take $\mu>0$ fixed, $m=l_{proj}-3R$ and $q=c_1l_{proj}\Delta E$ with
    $0<c_1<\mu/v$, and suppose the number $N$ of omitted terms meeting $K$
    grows subexponentially in $l_{proj}$, for instance polynomially when spheres
    in the graph have polynomial size. Then both terms of $\delta$ in
    Theorem~\ref{thm:hastings_truncated_facld} decay exponentially in
    $l_{proj}$, which is the form $JO(e^{-l_{proj}/l_{fac}})$ of
    eq.~(facld) \cite{gap:hastings05_truncated_dynamics}.
\end{remark}

\begin{theorem}[The exponential decay of eq.~(facld)]
    \label{thm:hastings_truncated_facld_decay}
    \lean{TNLean.Lattice.exists_norm_sub_starProjection_truncated_apply_le_exp,
        TNLean.Lattice.exists_norm_truncated_apply_sub_inner_smul_le_exp,
        Real.add_one_pow_mul_le_mul_exp}
    \leanok
    \uses{thm:hastings_truncated_facld}
    In the setting of Theorem~\ref{thm:hastings_truncated_facld}, let $J\ge0$,
    fix $\mu>0$ and $c_1>0$ with $c_1v<\mu$, and set
    $l_{fac}=2/\min(c_1\Delta E,\ \mu-c_1v)$. For each projection radius
    $l\in\mathbb{N}$ let $M_l$ be the truncated filter of $W$ at
    $q=c_1l\Delta E$, built from the terms of a set $s_l$ acting inside a
    region $K_l\supseteq S$, such that every omitted term meeting $K_l$ lies at
    distance at least $m=l-3R$ from $S$. Suppose the number $N_l$ of omitted
    terms meeting $K_l$ satisfies $N_l\le K_g(l+1)^k$. Then there is a constant
    $C_1'\ge0$, independent of $l$, such that for every $l$ and every ground
    vector $\psi$ the component of $M_l\psi$ orthogonal to the ground space
    has norm at most $\|W\|(3+C_1'J)e^{-l/l_{fac}}\|\psi\|$.
    If $\Omega$ is a unique ground state with gap $\Delta E$, there is also a
    constant $C_2'\ge0$, independent of $l$, such that for every $l$,
    $\|M_l\Omega-\langle\Omega,M_l\Omega\rangle\Omega\|\le
        \|W\|(3+C_2'J)e^{-l/l_{fac}}$. This is the form $JO(e^{-l_{proj}/l_{fac}})$
    of eq.~(facld) of \cite[lines 322--349]{Hastings2006Solving}, with the
    polynomial growth of the shell as an explicit hypothesis
    \cite{gap:hastings05_truncated_dynamics}.
\end{theorem}

\begin{proof}
    \leanok
    \uses{thm:hastings_truncated_facld}
    Write $\delta=\min(c_1\Delta E,\ \mu-c_1v)>0$. At $q=c_1l\Delta E$ one has
    $q/\Delta E=c_1l$, so the first term of the bound of
    Theorem~\ref{thm:hastings_truncated_facld} is
    $3\|W\|e^{-c_1\Delta El/2}\le3\|W\|e^{-\delta l/2}$. Since
    $\mu m\ge\mu l-3\mu R$, the exponential of the second term is at most
    $e^{3\mu R}e^{-(\mu-c_1v)l}\le e^{3\mu R}e^{-\delta l}$. For $a>0$ and
    $x\ge0$, the inequality $y^{k+1}/(k+1)!\le e^{y}$ at $y=a(x+1)$ gives
    $(x+1)^kx\le(k+1)!\,e^{a}a^{-(k+1)}e^{ax}$. With $a=\delta/2$ this bounds
    $N_l\,l$ by a constant times $e^{\delta l/2}$, so the second term is at
    most $\|W\|C_1'Je^{-\delta l/2}\|\psi\|$ in the ground-space estimate
    and $\|W\|C_2'Je^{-\delta l/2}$ in the unique-ground-state estimate.
\end{proof}
